\section{Modifications to chapter 12}

The following modifications shall be applied to chapter 12 (\emph{Overloading}
[\texttt{over}]): 

\begin{modification}[Modification of 12.2.2.2.2/3]
In unqualified function calls, the function is designated by an
\emph{id-expression} or
a \emph{splice-expression} \emph{E}. The set of candidate functions either is the set found by
name lookup (6.5) if \emph{E} is an \emph{id-expression} or is the set determined as specified
in 7.5.9 if \emph{E} is a \emph{splice-expression}. The set of candidate functions consists
either entirely of non-member functions or entirely of member functions of some
class \texttt{T}. In the former case or if \emph{E} is either a splice-expression or the address
of an overload set, the argument list is the same as the \emph{expression-list} in the
call. Otherwise, the argument list is the \emph{expression-list} in the call augmented
by the addition of an implied object argument as in a qualified function call.
If the current class is, or is derived from, \texttt{T}, and the keyword
\texttt{this} (7.5.3)
refers to it, \stdadd{the implied object argument is \texttt{(*this)}}.

\begin{itemize}[label=\textemdash]

\item
\stdremove{if the unqualified function call appears in a precondition assertion of a constructor or a postcondition
assertion of a destructor and overload resolution selects a non-static member function, the call is
ill-formed;}

\item
\stdremove{otherwise, the implied object argument is \texttt{(*this)}.}

\end{itemize}

Otherwise,

\begin{itemize}[label=\textemdash]

\item 
if overload resolution selects a non-static member function, the call is ill-formed;

\item
otherwise, a contrived object of type \texttt{T} becomes the implied object argument.

\end{itemize}

In the corresponding example \stdremove{remove the following lines}:
\begin{verbatim}
C()
  pre(a())          // error: implied this in constructor precondition
  pre(this->a())    // OK
  post(a());        // OK
~C()                
  pre(a())          // OK
  post(a())         // error: implied this in destructor postcondition
  post(this->a());  // OK
\end{verbatim}

\end{modification}
